\section{Where the analytic proof must be completed}
\label{sec:analytic}

This section records the genuine proof obligations behind the classical
criterion. It does not declare them axioms of $\ipra$ or count their
restatement as a proof. Robin's original argument is in
\cite[\S\S3--4]{Robin}; the Nicolas oscillation input is developed in
\cite{Nicolas1983,Nicolas2012}. A modern classical Lean reconstruction
provides a substantially more explicit dependency map \cite{Whidden}.

\subsection{An elementary extremal interpolation lemma}

Let $N,N'$ be integers with $3\le N<N'$. Suppose a real parameter
$\varepsilon>0$ and a constant $C$ give
\[
 \log A(j)\le C+\varepsilon\log j
       \quad(N\le j\le N'),
\]
with equality at both endpoints. Then the Robin logarithmic margin
\[
 M(j)=\log A(j)-\gamma-\log\log\log j
\]
satisfies
\begin{equation}\label{eq:interpolation}
 M(j)\le \max(M(N),M(N'))\qquad(N\le j\le N').
\end{equation}
Indeed, put $x=\log j$ and $f(x)=\varepsilon x-\log\log x$.
On $x>1$,
\[
 f''(x)=\frac{\log x+1}{x^2(\log x)^2}>0.
\]
Thus $f$ is convex and bounded above by the larger of its two endpoint
values on the interval. Adding $C-\gamma$ proves
\eqref{eq:interpolation}. One can replace the maximum statement by a
rational convex-combination upper bound when carrying explicit
certificates.

For colossally abundant endpoints the affine bound comes from maximizing
$A(j)j^{-\varepsilon}$. The displayed lemma itself only needs a
certified bound on the finite interval. Its proof does not require
deciding equality between unrelated transcendental threshold values.
A constructive implementation should allow ties or use overlapping
rational upper bounds, instead of assuming that every threshold
comparison is strictly decidable.

\subsection{The RH-dependent weighted explicit formula}

Put $H=\log n$,
\[
 S(H)=\frac1{\sqrt H\log H},
 \qquad c_0=\gamma+2-\log(4\pi).
\]
The sharp forward proof uses weighted Chebyshev errors
\[
 I_j(x)=\int_x^\infty(\psi(t)-t)t^{-j-1}
                      \frac{j\log t+1}{(\log t)^2}\,dt.
\]
Under RH a multiplicity-counted zero sum yields the estimates
\begin{align*}
 -B_j(x)-\frac{\log(2\pi)x^{-j}}{\log x}
 &\le I_j(x)\le B_j(x),\\
 B_j(x)&=c_0\frac{x^{1/2-j}}{\log x}
       \left(j+\frac1{\log x}
                  +\frac4{(2j-1)(\log x)^2}\right).
\end{align*}
The zero contribution is controlled by
$\sum_\rho|\rho|^{-2}=c_0$ under RH. These statements involve
infinite zero sums and integrals, not only evaluations of elementary
constants. The required arithmetic version must state finite contour
identities, prove their error estimates, control the zero multiplicities,
and provide bounds for discarded tails.

For example, a formal statement that a sequence of approximants converges
does not provide the modulus needed to replace this integral by a fixed
primitive-recursive observation. Likewise a type enumerating zeros is
not a proved bound on the tail of their reciprocal squares. The
quantitative packets already present in ProveIt are relevant candidate
arguments, but their stated proof status must be respected.

\subsection{A concrete sharp margin available after that work}

The classical reconstruction \cite{Whidden} organizes the forward estimate
as
\begin{equation}\label{eq:budget}
 M(n)\le E(H)-C(H)\qquad(H\ge20000,\ \RH),
\end{equation}
where
\begin{align*}
 C(H)&=\sum_{p\le H/2}
       \min\left(\frac{\log p}{H\log H},\frac1{p^2}\right),\\
 E(H)&=(2+c_0)\frac{H^{-1/2}}{\log H}
       +(c_0-2)\frac{H^{-1/2}}{(\log H)^2}\\
 &\quad +(8+4c_0)\frac{H^{-1/2}}{(\log H)^3}
       +2\frac{H^{-2/3}}{\log H}
       +\log(2\pi)\frac{H^{-1}}{\log H}.
\end{align*}
The first term is an explicit sum of losses forced by prime powers;
the second is an upper error estimate obtained from the weighted
explicit formula. The published reconstruction gives, for $H\ge74500$,
\begin{equation}\label{eq:rational-budgets}
 \frac{E(H)}{S(H)}
 \le\frac{3010451}{1332800}
 =\frac{113}{50}-\frac{1677}{1332800},
 \qquad
 \frac{C(H)}{S(H)}\ge\frac{113}{50}.
\end{equation}
These are recorded as concrete porting targets. This report has not
reproved their analytic premises in $\ipra$, nor independently rebuilt
the external Lean development.

Their implication for constructive strictness can, however, be derived
explicitly.

\begin{proposition}[Conditional positive separation]\label{prop:separation}
Assume the actual estimates \eqref{eq:budget} and
\eqref{eq:rational-budgets} for an integer $n$ with $\log n\ge74500$.
Let $a=1677/1332800$. Then
\[
 B(n)-A(n)\ge\frac{a}{2n^2}
             =\frac{1677}{2665600\,n^2}.
\]
Consequently the canonical interval test $\Pos(n,k)$ succeeds at
$k=2\ell(n)+12$.
\end{proposition}
\begin{proof}
The estimates imply $M(n)\le-aS(H)$.
For $n$ in this range, $\sqrt H\le n$ and $\log H\le n$, so
$S(H)\ge n^{-2}$ and $M(n)\le-a/n^2$.
Since $M(n)=\log(A(n)/B(n))$ and $B(n)\ge1$,
\[
 B(n)-A(n)\ge B(n)(1-e^{-a/n^2}).
\]
For $0\le u\le1$, $e^u\ge1+u$, and therefore
$1-e^{-u}\ge u/(1+u)\ge u/2$. Use $u=a/n^2$.

Finally $2^{-2\ell(n)-12}\le1/(4096n^2)<a/(2n^2)$.
An interval for $B(n)$ with width at most this quantity has lower
endpoint strictly above $A(n)$. The test \eqref{eq:pos} therefore
accepts at the stated primitive-recursive precision.
\end{proof}

This proposition shows why the strict inequality should not be treated
as a permanent obstruction. The forward argument already contains
room for a useful explicit positive witness. The unresolved question
is the justification of its analytic input in the chosen arithmetic
theory.

\subsection{The finite startup range is another real obligation}

The same external project reports three stages: exact factorizations
for $5041\le n<7560$, a common tangent up to $720720$, and $36$
rational logarithmic-height rows reaching the analytic range.
Its README reports ordinary kernel checks of the finite data and an
axiom audit using only \code{propext}, \code{Classical.choice}, and
\code{Quot.sound}. Those are reports of the external project, not
checks rerun in this task.

The finite verifier and its soundness proof are plausible objects to
port because the accepted rows are rational finite data. But a
single finite bound can be enormous, and a classical real-valued
soundness theorem still needs translation. The present archive's
verification through $100000$ does not cover this logarithmic-height
range and is not offered as a replacement for it.

A completed constructive argument should make the handoff using a
decidable integer cutoff, or overlapping rational height ranges.
It should not invoke the general real trichotomy
$\log n<74500$ or $\log n\ge74500$ without explaining its arithmetic
realization. The existing rational rows provide a useful approach to
that engineering problem, but no such port is claimed here.

\subsection{What the converse still needs}

The classical Nicolas--Landau mechanism supplies negative excursions
of $\log F(x)$ when an off-critical zero exists. In a convenient form,
there are $c>0$ and $0<b<1/2$ such that arbitrarily far out one has
\[
 -\log F(Q)\ge cQ^{-b}.
\]
Because $\vartheta(x)$ and the prime product are constant between
successive integers, replacing a real frontier by its integer part
only weakens the constant by a fixed factor. Thus an integer-frontier
form can be used.

The preceding LCM section proves the remaining finite implication
in full detail. Once an excursion beyond its explicit rational
cutoff is supplied, its output is an actual strict counterexample.
There is no additional number-theoretic asymptotic argument hidden
in that transfer.

The missing analytic proof is substantial. Its familiar ingredients
include analytic continuation of a weighted prime-function transform,
Landau's positive-transform principle, and the detection of a
singularity associated with an off-line zero. A claim of arbitrary
large excursions in classical analysis is not yet a
primitive-recursive witness bound in $\ipra$.

There are two defensible completion routes:
\begin{enumerate}[leftmargin=*]
\item Starting with a specified accepted zero certificate, prove a
quantitative transform theorem producing a frontier, a rational
excursion gap, and a primitive-recursive bound for their search.
The LCM transfer then gives \eqref{eq:R}.
\item Under $\Rstr$ and an accepted zero certificate, derive an
analytic contradiction inside a proved conservative arithmetic
framework. For the negative RH conclusion, a direct contradiction
can suffice; no positive counterexample selector need be asserted
unless it is actually obtained.
\end{enumerate}

The second route may be shorter. Neither is completed by citing
ordinary Lean's classical choice or by declaring the needed
oscillation theorem as a structure field.

\subsection{How conservation results could help, and their limit}

Sudac's work \cite{Sudac} develops real and complex analysis in
extensions conservative over PRA and obtains a PRA proof of the
prime number theorem. It is evidence that substantial analytic
number theory can admit such treatment. It is not a published
PRA proof of Robin's equivalence, and no implication from PNT
alone gives the sharp RH-sensitive estimates here.

A serious alternative to direct object-proof construction is to
check every needed analytic principle in an appropriate conservative
theory and formalize the finite witness statements there.
The hypothesis restrictions and the exact conservation class must
be verified. The transfer theorem in Section~\ref{sec:logic}
explains why a genuine classical PRA endpoint would suffice
afterward. This is a research route with a concrete endpoint,
not a substitute for the missing derivation.
